\documentclass{article}
\usepackage[margin=1in]{geometry}
\usepackage{amsmath,amssymb,amsthm}
\usepackage{hyperref}

\newtheorem{definition}{Definition}
\newtheorem{proposition}[definition]{Proposition}
\newtheorem{theorem}[definition]{Theorem}
\newtheorem{remark}[definition]{Remark}
\newcommand{\cA}{\mathcal A}
\newcommand{\cM}{\mathcal M}
\newcommand{\cR}{\mathcal R}
\newcommand{\cB}{\mathcal B}

\begin{document}

\section*{Blueprint 001: Tensor-product transport of traces}

\subsection*{Scope and correction}

The intended source statement is that a regional von Neumann algebra is
isomorphic to a tensor product of a type-I factor and a hyperfinite type
\(II_1\) factor.  Its original proof assumes rather than proves a compatible
infinite-coordinate map, product-state factorization, and transport to the
weak closures.  This blueprint therefore has two layers.  The first is a
complete algebraic theorem proved in Lean.  The second is the analytic model
identification, recorded with its necessary hypotheses but not represented by
an axiom in Lean.

\begin{definition}[Unital tracial functional, formal target TPT-001]
For a unital complex star algebra \(A\), a functional \(\tau:A\to\mathbb C\)
is unital and tracial when
\begin{equation*}
  \tau(1)=1,
  \qquad
  \tau(xy)=\tau(yx)\quad (x,y\in A).
\end{equation*}
It is algebraically positive when for every \(x\in A\),
\begin{equation*}
  \tau(x^*x)=r\in\mathbb R_{\geq0}\subset\mathbb C.
\end{equation*}
An \emph{algebraic tracial state} is a unital, algebraically positive,
tracial functional.  On a C\(^*\)-algebra this is the usual state condition,
except that normality is a separate topological property.
\end{definition}

\begin{theorem}[TraceTransportAlongStarIsomorphism, formal target TPT-002]
Let \(\Phi:A\to B\) be a unital \(*\)-isomorphism of unital complex star
algebras.  If \(\tau_B:B\to\mathbb C\) is an algebraic tracial state, then
\begin{equation*}
  \tau_A:=\tau_B\circ\Phi
\end{equation*}
is an algebraic tracial state on \(A\).
\end{theorem}

\begin{proof}
Since \(\Phi\) is unital,
\(\tau_A(1)=\tau_B(\Phi(1))=\tau_B(1)=1\).  Since \(\Phi\) is
multiplicative,
\begin{equation*}
 \tau_A(xy)=\tau_B(\Phi(x)\Phi(y))
 =\tau_B(\Phi(y)\Phi(x))=\tau_A(yx).
\end{equation*}
Finally,
\begin{equation*}
  \tau_A(x^*x)=\tau_B\bigl(\Phi(x)^*\Phi(x)\bigr)\in\mathbb R_{\geq0},
\end{equation*}
because \(\Phi\) preserves multiplication and the involution.  Lean checks
all three clauses of this algebraic statement.  Applying the same result to
\(\Phi^{-1}\) proves the converse: algebraic traciality is reflected as well
as preserved by a star-algebra equivalence.
\end{proof}

\begin{proposition}[FiniteMatrixNormalizedTrace, formal target FMT-001]
For a nonempty finite index set \(I\), the functional
\begin{equation*}
  \operatorname{tr}_I(a)
  :=\frac{1}{|I|}\operatorname{Tr}(a),
  \qquad a\in M_I(\mathbb C),
\end{equation*}
is an algebraic tracial state.
\end{proposition}

\begin{proof}
The matrix trace of the identity is \(|I|\), so the normalization gives
\(\operatorname{tr}_I(1)=1\).  Cyclicity follows from
\(\operatorname{Tr}(ab)=\operatorname{Tr}(ba)\).  For positivity, the Lean
proof computes
\begin{equation*}
  \operatorname{Tr}(a^*a)=\sum_{i,j\in I}\lvert a_{ji}\rvert^2\geq0.
\end{equation*}
Thus the Lean theorem
\texttt{normalizedMatrixTrace\_isAlgebraicTracialState} proves that the
concrete normalized matrix trace is an algebraic tracial state.
\end{proof}

\begin{proposition}[FiniteSupportCornerTrace, formal target FSC-001]
Let \(q\cA q\) be a finite support corner and assume an explicit unital
\(*\)-isomorphism
\begin{equation*}
  e:q\cA q\longrightarrow M_I(\mathbb C)
\end{equation*}
for a nonempty finite set \(I\).  Then
\begin{equation*}
  \tau_q:=\operatorname{tr}_I\circ e
\end{equation*}
is an algebraic tracial state on \(q\cA q\).
\end{proposition}

\begin{proof}
Apply Theorem TPT-002 to the normalized matrix trace from Proposition
FMT-001.  The Lean theorem
\texttt{finiteCornerTrace\_isAlgebraicTracialState}
checks this transport.  Identifying an actual stabilizer support corner with a
full matrix factor, and proving that its state equals this normalized trace,
remain separate stabilizer-theoretic obligations.
\end{proof}

\begin{proposition}[FiniteTensorProductTrace, formal target FTT-001]
For nonempty finite index sets \(P\) and \(R\), let
\begin{equation*}
  M_P(\mathbb C)\otimes_{\mathrm{alg}}M_R(\mathbb C)
  \cong M_{P\times R}(\mathbb C)
\end{equation*}
be the Kronecker \(*\)-isomorphism.  Pulling the normalized trace on
\(M_{P\times R}(\mathbb C)\) back along this isomorphism gives an algebraic
tracial state \(\operatorname{tr}_{P\otimes R}\), and for elementary tensors,
\begin{equation*}
  \operatorname{tr}_{P\otimes R}(a\otimes b)
  =\operatorname{tr}_P(a)\operatorname{tr}_R(b).
\end{equation*}
\end{proposition}

\begin{proof}
The Kronecker map preserves multiplication and the involution.  The target
matrix trace is an algebraic tracial state by Proposition FMT-001, so
Theorem TPT-002 transports it to the algebraic tensor product.  On elementary
tensors, the identity follows from
\begin{equation*}
  \operatorname{Tr}(a\otimes b)=\operatorname{Tr}(a)\operatorname{Tr}(b)
\end{equation*}
and \(|P\times R|=|P||R|\).  The Lean theorem
\texttt{finiteTensorTrace\_tmul} verifies this formula using mathlib's
Kronecker star-algebra equivalence and trace calculation.
Its checked append-left and append-right lemmas prove compatibility of
normalized traces with adjoining a finite factor by tensoring with the
identity.
\end{proof}

\begin{proposition}[Finite type-I amplification, formal target FAM-001]
Let \(R\) be a unital complex star algebra with algebraic tracial state
\(\tau_R\), and let \(I\) be a nonempty finite index set.  The normalized
diagonal functional on the matrix amplification is
\begin{equation*}
  \tau_{I,R}(a):=\frac{1}{|I|}\sum_{i\in I}\tau_R(a_{ii}),
  \qquad a\in M_I(R).
\end{equation*}
It is an algebraic tracial state.  Consequently, if a region algebra \(A\)
is supplied with an explicit complex \(*\)-isomorphism
\begin{equation*}
  e:A\longrightarrow M_I(R),
\end{equation*}
then \(\tau_{I,R}\circ e\) is an algebraic tracial state on \(A\).
\end{proposition}

\begin{proof}
Unitality follows by evaluating the diagonal of the identity.  For cyclicity,
expand both diagonal traces as a double finite sum and use cyclicity of
\(\tau_R\) term by term, then exchange the two finite sums.  For positivity,
the Lean proof establishes
\begin{equation*}
  \sum_i \tau_R\bigl((a^*a)_{ii}\bigr)
  =\sum_{i,j}\tau_R(a_{ji}^*a_{ji}),
\end{equation*}
which is a finite sum of nonnegative real numbers after applying positivity of
\(\tau_R\).  The Lean theorem
\texttt{normalizedMatrixTraceOver} amplification result checks these steps;
the corresponding induced-trace result checks the final transport through
\(e\).  In the intended factor model,
\(M_I(R)\) is the algebraic finite type-I amplification of the tracial
factor; identifying it with an analytic von Neumann tensor product remains a
separate completion theorem.
\end{proof}

\begin{proposition}[Canonical finite type-I tensor identification, formal target TMT-001]
\label{prop:canonical_matrix_tensor}
Let $R$ be a unital complex star algebra whose complex scalar action
commutes with the involution, and let $I$ be finite.  There is a canonical
complex star-algebra isomorphism
\begin{equation*}
  M_I(\mathbb C)\otimes_{\mathrm{alg}}R
  \cong M_I(R).
\end{equation*}
If $\tau_R$ is an algebraic tracial state and $I$ is nonempty, the
pullback of the normalized diagonal trace on $M_I(R)$ is an algebraic
tracial state.  For $a\in M_I(\mathbb C)$ and $r\in R$, it obeys
\begin{equation*}
  \tau_{I\otimes R}(a\otimes r)=\operatorname{tr}_I(a)\tau_R(r).
\end{equation*}
\end{proposition}

\begin{proof}
The Lean map first exchanges the two tensor factors and then uses mathlib's
canonical matrix--tensor algebra equivalence.  Its involution preservation is
proved entrywise, including the use of the central complex scalar action in
$R$.  The proof of the displayed trace identity expands the normalized
diagonal trace and uses cyclicity of $\tau_R$ to commute $r$ past each
complex scalar.  The checked map is
\texttt{matrixTensorStarAlgEquiv}; the state and elementary-tensor trace
theorems are verified in the same \texttt{TensorMatrixEquiv} module.
\end{proof}

\begin{proposition}[Concrete algebraic dyadic trace system, formal target DTT-001]
\label{prop:concrete_dyadic_trace}
Put $A_0=\mathbb C$ and recursively define
\begin{equation*}
  A_{n+1}=M_2(A_n),
  \qquad
  \iota_n(a)=\operatorname{diag}(a,a).
\end{equation*}
The preceding canonical finite type-I identification interprets this as the
usual operation $a\mapsto 1_2\otimes a$.  Define
\begin{equation*}
  \tau_0(z)=z,
  \qquad
  \tau_{n+1}(b)=\frac12\sum_{i=0}^1\tau_n(b_{ii}).
\end{equation*}
Then every $\tau_n$ is an algebraic tracial state and, for every iterated
inclusion $\iota_{n,m}:A_n\to A_m$,
\begin{equation*}
  \tau_m\bigl(\iota_{n,m}(a)\bigr)=\tau_n(a).
\end{equation*}
Consequently the algebraic directed limit
$\mathcal U_{\mathrm{alg}}:=\varinjlim_n A_n$ has a canonical algebraic tracial
state whose composition with the canonical map from $A_n$ is $\tau_n$.
\end{proposition}

\begin{proof}
The base trace is the identity on \(\mathbb C\), whose value on a star-square
is the complex norm-square.  The successor trace is the finite matrix
amplification from Proposition FAM-001.  The diagonal embedding contributes
the same coefficient algebra element to both diagonal entries, so its
normalized trace is unchanged.  Lean then proves compatibility for all
iterated inclusions by induction and applies Theorem AIT-002 to the resulting
directed system.  The checked endpoint is
\texttt{qubitUHFTrace\_isAlgebraicTracialState}, with
\texttt{qubitUHFTrace\_of} recording its finite-stage formula.  This is
not yet the norm-completed UHF C$^*$-algebra or the weak closure giving the
hyperfinite type $II_1$ factor.  The finite stages now also carry explicit
ordered Pauli-word star systems and Pauli-word bases.  Lean verifies that the
identity-qubit embeddings preserve their words, basis vectors, and both
multiplier functions under coherent label maps.
\end{proof}

\begin{proposition}[Finite pure factor times the concrete tracial product,
formal target FPU-001]
\label{prop:finite_pure_qubit_product}
Let $I$ be a nonempty finite set and let
$\mathcal U_{\mathrm{alg}}$ be the algebraic dyadic direct limit from
Proposition DTT-001, equipped with its checked trace
$\tau_{\mathcal U}$.  Then
\begin{equation*}
  M_I(\mathbb C)\otimes_{\mathrm{alg}}\mathcal U_{\mathrm{alg}}
\end{equation*}
carries an algebraic tracial state $\tau_{I,\mathcal U}$ satisfying
\begin{equation*}
  \tau_{I,\mathcal U}(a\otimes u)
  =\operatorname{tr}_I(a)\tau_{\mathcal U}(u).
\end{equation*}
\end{proposition}

\begin{proof}
This is the concrete instantiation of Proposition TMT-001 with
$R=\mathcal U_{\mathrm{alg}}$.  The preceding direct-limit
star-module theorem supplies the scalar--involution compatibility needed by
that tensor construction.  Lean checks the resulting state and
elementary-tensor formula in the finite-pure-product module.  It remains an
algebraic product; it does not assert a normal trace on a von Neumann tensor
product.
\end{proof}

\begin{proposition}[Finite pure/mixed coordinate product, formal target PMS-001]
\label{prop:finite_pure_mixed_coordinate_product}
Fix a sequence $p:\mathbb N\to\{0,1\}$.  On the finite qubit tower, define
$\Omega_{p,0}$ to be the scalar state and, after adjoining one qubit on the
left, define
\begin{equation*}
  \Omega_{p,n+1}(a\otimes x)=
  \begin{cases}
    \langle 0|a|0\rangle\,\Omega_{p,n}(x),&p(n)=1,\\
    \operatorname{tr}_2(a)\,\Omega_{p,n}(x),&p(n)=0.
  \end{cases}
\end{equation*}
Here the displayed tensor notation is transported through the canonical
matrix--tensor star equivalence.  Then every $\Omega_{p,n}$ is an algebraic
state, the family is compatible with identity-qubit embeddings, and it induces
an algebraic state on the algebraic qubit direct limit.  When $p=0$, this
state is the normalized algebraic UHF trace.
\end{proposition}

\begin{proof}
For a pure new coordinate, positivity is the finite sum of the coefficient
state values on the entries of the first column of a matrix star-square.  For a
mixed coordinate, it is the normalized matrix amplification of the coefficient
state.  Both functionals evaluate a scalar identity-qubit embedding as the
coefficient state, which proves compatibility by induction.  Lean checks these
claims and the direct-limit state in
\texttt{QubitMatrixState.lean} and
\texttt{FinitePureMixedQubitState.lean}.  A pure vector state is not claimed
to be tracial; only the all-mixed specialization is the tracial UHF state.  On
an ordered Pauli word, the state is one precisely when every pure coordinate
has zero X exponent and every mixed coordinate has zero X/Z label; it is zero
otherwise.  Lean separately verifies the expected obstruction already for one
pure qubit, and now for a pure coordinate at an arbitrary finite or
direct-limit stage: its vector state is not cyclic, and therefore is not
tracial.  In fact, Lean proves that the algebraic UHF product state is tracial
if and only if every coordinate is mixed, in which case it is the normalized
algebraic UHF trace.
\end{proof}

\begin{proposition}[Coherent phase-corrected coordinate states, formal target
CPS-001]
\label{prop:coherent_phase_coordinate_state}
Let $(A_n,\alpha_{n,m})$ be a directed system of complex star algebras, and
let $b_n:V_n\to A_n$ be a complex linear basis at every finite stage.  Suppose
there are star-algebra equivalences
\begin{equation*}
  e_n:A_n\longrightarrow \mathcal U_n
\end{equation*}
to the finite qubit stages, label maps $t_n:V_n\to\mathbb F_2^{2n}$, and
complex phases $\lambda_n:V_n\to\mathbb C$ such that
\begin{equation*}
  e_n(b_n(v))=\lambda_n(v)W_n(t_n(v)).
\end{equation*}
If the source basis embeddings, phases, and labels commute with every
enlargement of the finite stage, then the $e_n$ commute with the source and
target connecting maps on every element of $A_n$.  They therefore induce a
star-algebra equivalence from the algebraic source limit to the algebraic
qubit-UHF limit.  Hence, for every pure/mixed coordinate pattern $p$, the
qubit product state pulls back to an algebraic direct-limit state $\omega_p$
satisfying
\begin{equation*}
  \omega_p(\iota_n(b_n(v)))
  =\lambda_n(v)\,\mathbf 1_{\mathrm{sel}(p,n)}(t_n(v)),
\end{equation*}
where the indicator is one exactly when pure coordinates have zero X exponent
and mixed coordinates have the identity Pauli label.  When every coordinate is
mixed, this construction pulls the algebraic UHF trace back to an algebraic
tracial state on the source limit.  More generally, its pure/mixed pullback is
algebraically tracial if and only if every coordinate is mixed.
\end{proposition}

\begin{proof}
The two routes around the finite-stage square are complex-linear maps out of
$A_n$.  They agree on the basis $b_n(V_n)$ by the three stated coherence
conditions, hence agree everywhere by basis extensionality.  The compatible
finite-stage equivalences then pull the checked qubit product-state family
back to the algebraic source limit.  Evaluating a basis word first uses its
phase-corrected Pauli image, then the finite Pauli-word selection rule from
Proposition~\ref{prop:finite_pure_mixed_coordinate_product}.  Lean checks the
basis-extensionality step in \texttt{CoherentBasisEquiv.lean} and the stated
state construction and formula in \texttt{CoherentPhaseQubitState.lean}.
The exact recursive-to-coordinate identification is checked in
\texttt{FiniteQubitPauliSelection.lean}; it is the selected-Z condition used
in the displayed formula.
With the additional finite-stage hypothesis that these selected target labels
are exactly the actual source stabilizer basis words,
\texttt{CoherentPhaseQubitState.lean} derives the same formula with the
indicator of the actual stabilizer predicate.  Thus the only remaining input
at this point is the regional membership identification itself.
This proposition does not construct the regional bases, phases, or coherent
label maps; it makes those finite geometric conditions the exact remaining
input.
\end{proof}

\begin{proposition}[Pauli-word phase-compatibility criterion, formal target
PCC-001]
\label{prop:pauli_cocycle_compatibility}
Let $p_V(v)$ and $p_W(w)$ be two Pauli-word conventions, with multiplication
and involution phases defined by
\begin{align*}
  p_V(v)p_V(v')&=m_V(v,v')p_V(v+v'),
  &p_V(v)^*&=s_V(v)p_V(-v),\\
  p_W(w)p_W(w')&=m_W(w,w')p_W(w+w'),
  &p_W(w)^*&=s_W(w)p_W(-w).
\end{align*}
For an additive coordinate map $T:V\to W$ and a phase function
$\lambda:V\to\mathbb C$, the phase-corrected assignment
\begin{equation*}
  p_V(v)\longmapsto \lambda(v)p_W(Tv)
\end{equation*}
respects both displayed word relations provided
\begin{align*}
  \lambda(v)\lambda(v')m_W(Tv,Tv')
    &=m_V(v,v')\lambda(v+v'),\\
  \overline{\lambda(v)}s_W(Tv)
    &=s_V(v)\lambda(-v).
\end{align*}
\end{proposition}

\begin{proof}
Expand the product of the two target words using the target multiplication
relation.  The first phase equation changes its scalar coefficient to the
source coefficient multiplying the corrected word at $v+v'$.  Similarly,
take the star of the corrected word, use the target star relation, and apply
the second phase equation.  Lean verifies these two word-level calculations
in \texttt{PauliCocycleCompatibility.lean}.  The proposition does not claim
that symplecticity yields $\lambda$, nor does it extend the assignment to a
star-algebra equivalence: those require respectively a phase-correction
construction and a Pauli-word presentation or linear-independence theorem.
\end{proof}

\begin{proposition}[Algebraic direct-limit transport, formal target ALC-001]
\label{prop:algebraic_direct_limit_transport}
Let $(A_i,\alpha_{ij})$ and $(B_i,\beta_{ij})$ be directed systems of complex
star algebras.  If complex star-algebra homomorphisms
$\phi_i:A_i\to B_i$ obey
\begin{equation*}
  \phi_j\circ\alpha_{ij}=\beta_{ij}\circ\phi_i
  \qquad (i\leq j),
\end{equation*}
then they induce a complex star-algebra homomorphism
\begin{equation*}
  \varinjlim_i A_i\longrightarrow\varinjlim_i B_i.
\end{equation*}
If every $\phi_i$ is a star-algebra equivalence, the induced map is a
star-algebra equivalence.
\end{proposition}

\begin{proof}
The map sends the class of a stage element $a\in A_i$ to the class of
$\phi_i(a)\in B_i$.  The displayed compatibility makes this independent of
the representative.  Addition, multiplication, scalar action, and the star
can all be computed at one common stage.  For the equivalence assertion,
apply the same construction to the compatible inverses and verify the two
composites on stage representatives.  Lean verifies both claims in
\texttt{AlgebraicLimitStarHom.lean}.  This is algebraic only: a C$^*$-map
requires the additional proof that the finite-stage maps are isometric and
compatible with the chosen norm completions.
\end{proof}

\begin{proposition}[C-star continuity of a supplied coordinate equivalence,
formal target CSE-001]
Let $A$ and $B$ be complex C$^*$-algebras, and let
\begin{equation*}
  e:A\longrightarrow B
\end{equation*}
be a complex star-algebra equivalence.  Then $e$ is an isometry and is
continuous with continuous inverse.  Moreover, if a complex-linear functional
$\tau:B\to\mathbb C$ is continuous, then so is its pullback
\begin{equation*}
  \tau\circ e:A\longrightarrow\mathbb C.
\end{equation*}
\end{proposition}

\begin{proof}
For a complex star-algebra equivalence, norm contractivity applies to $e$ and
$e^{-1}$, giving equality of norms.  Thus both maps are continuous.  The
final assertion is closure of continuous maps under composition.  Lean checks
these statements in \texttt{CStarEquivExtension.lean}.  This proposition is
conditional on actual C$^*$-algebra structures: it does not construct the
norms or completions of the algebraic regional and coordinate limits.
\end{proof}

\begin{proposition}[Extension from a Pauli-word basis, formal target WBE-001]
\label{prop:word_basis_extension}
Let $A$ and $B$ be complex algebras, let $(p_i)_{i\in I}$ be a complex linear
basis of $A$, and let $L:A\to B$ be complex linear.  If
\begin{equation*}
  L(p_i p_j)=L(p_i)L(p_j)
  \qquad (i,j\in I),
\end{equation*}
then $L(xy)=L(x)L(y)$ for all $x,y\in A$.  If moreover $L(1)=1$, it is a
complex algebra homomorphism.  If $A$ and $B$ are star algebras,
\begin{equation*}
  L(p_i^*)=L(p_i)^*
  \qquad (i\in I),
\end{equation*}
then $L(x^*)=L(x)^*$ for every $x\in A$.  Consequently, if $L$ is bijective,
it is a complex star-algebra equivalence.
\end{proposition}

\begin{proof}
For fixed $y$, both maps
\begin{equation*}
  x\longmapsto L(xy),
  \qquad
  x\longmapsto L(x)L(y)
\end{equation*}
are complex linear.  It is therefore enough to compare them on each $p_i$.
For fixed $i$, the same argument in the $y$ variable reduces that comparison
to the assumed equality on $(p_i,p_j)$.  The Lean theorem
\texttt{LinearMap.map\_mul\_of\_basis} formalizes this twofold basis
argument, and \texttt{algHomOfBasis} adds the unitality clause.  This result
also treats the star condition: the set of $x$ satisfying
$L(x^*)=L(x)^*$ is a complex submodule, because both sides are
conjugate-linear in $x$, and it contains the basis.  The checked definition
\texttt{starAlgEquivOfBasis} combines the unital multiplicativity theorem,
this star argument, and bijectivity into a star-algebra equivalence.
\end{proof}

\begin{proposition}[Explicit phase correction (formal target QPR-001)]
\label{prop:explicit_pauli_phase_correction}
Let $T:V\to W$ be an isomorphism of $\mathbb F_2$-vector spaces preserving
the alternating commutator forms.  Suppose the two Pauli-word conventions
have multipliers
\begin{equation*}
  m_V(v,v')=(-1)^{B_V(v,v')},
  \qquad
  m_W(w,w')=(-1)^{B_W(w,w')},
\end{equation*}
where $B_V$ and $B_W$ are $\mathbb F_2$-bilinear and their symmetrizations are
the respective commutator forms, and suppose the involution conventions are
$s_V(v)=(-1)^{B_V(v,v)}$ and
$s_W(w)=(-1)^{B_W(w,w)}$.  Choose an ordered basis $(e_i)_{i\in I}$ of $V$.
Then a phase function $\lambda:V\to\{1,i,-1,-i\}$ satisfying
Proposition~\ref{prop:pauli_cocycle_compatibility} exists.
\end{proposition}

\begin{proof}
Put
\begin{equation*}
  D(v,v')=B_V(v,v')+B_W(Tv,Tv')\in\mathbb F_2.
\end{equation*}
The symplectic assumption says that $D$ is symmetric.  Write a vector as the
finite sum $v=\sum_i a_i e_i$, with $a_i\in\mathbb F_2$, and define
\begin{equation*}
  q(v)=\sum_i a_iD(e_i,e_i)
       +2\sum_{i<j}a_i a_jD(e_i,e_j)\pmod 4.
\end{equation*}
Here each $\mathbb F_2$ coefficient is represented by $0$ or $1$ before the
sum is read in $\mathbb Z/4\mathbb Z$.  Both sums are finite, even when $I$ is
infinite.  Expanding the two sums over $\mathbb F_2$ gives
\begin{equation*}
  q(v+v')=q(v)+q(v')+2D(v,v')\pmod 4.
\end{equation*}
Set $\lambda(v)=i^{q(v)}$.  This identity gives
\begin{equation*}
  \lambda(v)\lambda(v')m_W(Tv,Tv')
  =m_V(v,v')\lambda(v+v'),
\end{equation*}
which is the multiplication equation.  Also $q(v)$ modulo $2$ is
$D(v,v)$, because the second sum has coefficient $2$.  The stated involution
conventions, this parity identity, and $-v=v$ give
\begin{equation*}
  \overline{\lambda(v)}s_W(Tv)
  =s_V(v)\lambda(v)=s_V(v)\lambda(-v),
\end{equation*}
which is the star equation in
Proposition~\ref{prop:pauli_cocycle_compatibility}.  The Lean module
\texttt{QuadraticPhaseRefinement.lean} formalizes both implications from the
polarization and diagonal identities to the respective compatibility
equations.  Thus
$p_V(v)\mapsto\lambda(v)p_W(Tv)$ respects the prescribed multiplication and
star relations on Pauli-word generators.  The finite-support formula for $q$
is now defined in Lean and checked on zero, singleton, and two-coordinate
words, together with the recurrence obtained by adjoining a largest support
coordinate.  The associated finite double-sum discrepancy pairing is checked
to be additive in both arguments and symmetric when $D$ is symmetric; the
global polarization and diagonal identities are now formalized.  Once the
actual source and target Pauli exponents are identified with this pairing and
a phase-corrected word-basis linear equivalence is supplied, the checked
definition \texttt{coordinateQuadraticPhaseStarAlgEquivOfBasis} produces a
star-algebra equivalence between the source and target word algebras.  When
the target is the concrete finite-pure matrix algebra tensored with the
algebraic qubit-UHF direct limit, the checked theorem
\texttt{exists\_coordinateQuadraticPhase\_finitePureQubitUHFTrace} also pulls
the product trace back along this equivalence and proves its elementary-tensor
formula.  This remains an algebraic statement; the C$^*$ and von Neumann
extensions are not part of the theorem.
\end{proof}

\begin{proposition}[Calibration by the stabilizer character, formal target SCC-001]
\label{prop:stabilizer_character_calibration}
Let $W\subset V$ be the label subspace of the stabilizer group, and let
$P(v)P(w)=\sigma(v,w)P(v+w)$ be a fixed Pauli section.  The actual stabilizer
operators have the form
\begin{equation*}
  S(v)=\chi(v)P(v)\qquad (v\in W),
\end{equation*}
where their group law is equivalent to
\begin{equation*}
  \chi(v)\chi(w)\sigma(v,w)=\chi(v+w)
  \qquad (v,w\in W).
\end{equation*}
Let $V_{\mathrm{coord}}$ be the coordinate-label space.  Suppose
$T:V_{\mathrm{coord}}\to V$ is the adapted symplectic coordinate isomorphism
and $\lambda_0:V_{\mathrm{coord}}\to U(1)$ is a phase correction from
Proposition~\ref{prop:explicit_pauli_phase_correction}.  Then there is a
phase correction $\lambda:V_{\mathrm{coord}}\to U(1)$, obtained from
$\lambda_0$ by a character calibration, such that
\begin{equation*}
  \lambda(0,b)P(T(0,b))=S(T(0,b))
  \qquad\text{for every coordinate }(0,b)\in T^{-1}(W).
\end{equation*}
\end{proposition}

\begin{proof}
The quotient
\begin{equation*}
  \eta(v):=\chi(Tv)\lambda_0(v)^{-1}\qquad (v\in T^{-1}W)
\end{equation*}
is an ordinary character of the $\mathbb F_2$-vector space $T^{-1}W$:
the two cocycle relations cancel.  A character of a subspace extends to the
ambient vector space after extending a basis of the subspace to a basis of
$V_{\mathrm{coord}}$.  Multiply $\lambda_0$ by such an extension of $\eta$.  Ordinary
characters of an $\mathbb F_2$-vector space take values in $\{1,-1\}$, so they
preserve the multiplier and star compatibility equations, while
the displayed equality holds by construction.  Thus the corrected coordinate
$Z$-words are the actual stabilizers, not merely phase-ambiguous Pauli labels.
The quotient-cancellation, subspace-extension, finite-field involution, and
both Pauli-word compatibility-preservation steps are formalized in
\texttt{StabilizerCharacterCalibration.lean}.  The initial finite-support
quadratic phase correction is now formalized; its identification with the
actual regional Pauli convention remains a separate task.
\end{proof}

\begin{proposition}[Adapted-coordinate stabilizer membership, formal target
ASC-001]
\label{prop:adapted_coordinate_stabilizer_membership}
Let $V$ have a Hamel basis $\{z_i,x_i:i\in I\}$ over $\mathbb F_2$, let
$P\subseteq I$, and put
\begin{equation*}
  W_P=\operatorname{span}_{\mathbb F_2}\{z_i:i\in P\}.
\end{equation*}
Write a finite coordinate word as
\begin{equation*}
  v=\sum_{i\in I} a_i x_i+\sum_{i\in I}b_i z_i.
\end{equation*}
Then
\begin{equation*}
  v\in W_P
  \quad\Longleftrightarrow\quad
  a_i=0\ \text{for every }i,\qquad
  b_i=0\ \text{for every }i\notin P.
\end{equation*}
\end{proposition}

\begin{proof}
The coordinate equivalence identifies the selected-Z span with the finitely
supported functions on the subset $\{\mathrm{inl}(i):i\in P\}$ of the
disjoint X/Z coordinate index.  Membership is therefore precisely support
containment in that subset, which is equivalent to the displayed pointwise
conditions.  Lean proves both the equality with the selected-Z span and the
membership equivalence in \texttt{AdaptedStabilizerSubspace.lean}.  This
settles the linear selection criterion required in the density-matrix step.

\texttt{FiniteQubitPauliSelection.lean} now maps the recursive finite-qubit
labels into this finitely supported Z/X coordinate space and proves that its
pure/mixed selection predicate is precisely the displayed membership condition.
Identifying the concrete regional stabilizer subspace with this selected
coordinate span and matching its Pauli phases remain separate hypotheses.
\end{proof}

\begin{theorem}[AlgebraicInductiveLimitTrace, formal target AIT-001]
Let \((A_i,\alpha_{ij})\) be a directed system of unital complex algebras
with unital algebra homomorphisms.  Suppose each \(A_i\) carries a unital
cyclic linear functional \(\tau_i\), and that
\begin{equation*}
  \tau_j\bigl(\alpha_{ij}(a)\bigr)=\tau_i(a)
  \qquad (i\leq j, a\in A_i).
\end{equation*}
Then the algebraic direct limit carries a unique induced complex-linear
functional \(\tau_{\mathrm{alg}}\) with
\begin{equation*}
  \tau_{\mathrm{alg}}\circ\iota_i=\tau_i,
\end{equation*}
and \(\tau_{\mathrm{alg}}\) is unital and cyclic.
\end{theorem}

\begin{proof}
The compatibility equation makes the assignment
\(\iota_i(a)\mapsto\tau_i(a)\) independent of the representative in the
quotient defining the directed algebraic limit.  If two limit elements are
represented at one common stage, multiplication and cyclicity can be computed
there, proving cyclicity in the limit.  The Lean theorem
\texttt{algebraicLimitTrace\_isUnitalCyclic} proves this quotient-level
argument using mathlib's directed algebraic direct limit.
\end{proof}

\begin{theorem}[AlgebraicInductiveLimitTracialState, formal target AIT-002]
Let \((A_i,\alpha_{ij})\) be as in Theorem AIT-001, now with unital complex
star algebras and star-preserving connecting maps.  If every \(\tau_i\) is an
algebraic tracial state and the same compatibility condition holds, then the
algebraic direct limit has a canonical involution and its induced functional
is unital, cyclic, and positive on star-squares:
\begin{equation*}
  \tau_{\mathrm{alg}}(z^*z)\in\mathbb R_{\geq0}
  \qquad (z\in\varinjlim_i A_i).
\end{equation*}
\end{theorem}

\begin{proof}
Define the star of the class of \((i,a)\) to be the class of \((i,a^*)\).
This is well-defined because each connecting map preserves the involution.
For a representative \(z=\iota_i(a)\), both \(z^*z\) and its trace are
represented at the same stage by \(a^*a\), where positivity is an assumption.
The Lean module \texttt{AlgebraicStarDirectLimit} checks the quotient-level
well-definedness of the involution, its involutivity, additivity, and reversal
of multiplication, as well as the displayed positivity conclusion.  This is
still a statement about the \emph{algebraic} direct limit: it asserts neither
continuity nor normality after a completion.
\end{proof}

\begin{theorem}[Algebraic quasi-local trace transport, formal target ALT-001]
Let \((A_i,\alpha_{ij})\) and \((D_i,\delta_{ij})\) be directed systems of
unital complex star algebras.  Suppose star-algebra equivalences
\(e_i:A_i\to D_i\) satisfy
\begin{equation*}
  e_j\circ\alpha_{ij}=\delta_{ij}\circ e_i,
\end{equation*}
and compatible target algebraic tracial states \(\tau_i\) are given.  Then the
induced algebraic direct-limit star-algebra equivalence pulls the target-limit
trace back to an algebraic tracial state on \(\varinjlim_i A_i\), with
\begin{equation*}
  \tau_{\mathrm{source}}(\iota_i(a))=\tau_i(e_i(a)).
\end{equation*}
\end{theorem}

\begin{proof}
The finite-stage equivalences induce the algebraic-limit star-algebra
equivalence, and the compatible target-stage states glue by Theorem AIT-002.
Pull back the latter state.  The finite-stage formula follows by evaluating
both constructions on a representative.  Lean checks this argument in
\texttt{AlgebraicLimitTrace\-Transport.lean}; it is algebraic and does not
assert a C$^*$-continuous or normal extension.
\end{proof}

\begin{theorem}[Conditional finite-pure-sector trace]
Assume a normal \(*\)-isomorphism
\begin{equation*}
  \Phi:\cM_\Gamma \longrightarrow
  \cB(\mathcal K)\,\bar\otimes\,\cR,
\end{equation*}
where \(\dim\mathcal K<\infty\), \(\cR\) is the hyperfinite type
\(II_1\) factor, and the tensor product carries the product of the normalized
matrix trace and the unique normalized trace of \(\cR\).  Then its pullback
along \(\Phi\) is a faithful normal tracial state on \(\cM_\Gamma\).
\end{theorem}

\begin{proof}
The finite-dimensional normalized matrix trace and the normalized trace of
\(\cR\) are faithful normal tracial states.  Their normal tensor product is a
faithful normal tracial state; transport it by the normal isomorphism.  The
algebraic tracial part is Theorem TPT-002.  Normality and faithfulness are
preserved by a normal \(*\)-isomorphism.
\end{proof}

\begin{remark}[Infinite pure sector]
If \(\mathcal K\) is infinite dimensional, \(\cB(\mathcal K)\) has no
tracial state.  Hence \(\cB(\mathcal K)\bar\otimes\cR\) cannot acquire a
tracial state by tensoring the trace on \(\cR\) with a trace on the first
factor.  The correct object in this case is the faithful normal semifinite
trace obtained from the standard semifinite trace on \(\cB(\mathcal K)\) and
the finite trace on \(\cR\).  This is the type \(II_\infty\) case, not the
type \(II_1\) case.  The Lean module
\texttt{ProperlyInfiniteTraceObstruction.lean} checks the algebraic core: if
two isometries have range projections summing to the identity, unital cyclicity
would give \(1=1+1\) for the trace value of the identity.  Identifying such
isometries in a particular infinite von Neumann representation is a separate
analytic step.
\end{remark}

\subsection*{Corrected analytic target}

The analytic theorem needed to connect the checked algebraic results to a
stabilizer region should be stated with the following explicit hypotheses.
Let \(I\) be the coordinate set, let \(P\subset I\) label the pure
coordinates, and put \(C=I\setminus P\).  Assume:
\begin{enumerate}
\item \(V_{\mathcal G}\) has an adapted Hamel symplectic basis
\(\{z_i,x_i:i\in I\}\), every basis vector has finite lattice support, and
\(W_{\mathcal G}=\operatorname{span}_{\mathbb F_2}\{z_i:i\in P\}\);
\item the symplectic coordinate map is equipped with the finite-support
Pauli phase convention of Proposition~\ref{prop:explicit_pauli_phase_correction},
\item the phase convention is calibrated to the stabilizer character as in
Proposition~\ref{prop:stabilizer_character_calibration},
\item its induced map on finite Pauli words extends to a C\(^*\)-isomorphism
from the coordinate UHF algebra onto the regional quasi-local algebra; and
\item the restriction of the global GNS representation to that regional
algebra is quasi-equivalent to the GNS representation of the restricted
product state.
\end{enumerate}
The first condition uses finite support of each coordinate pair, not a
uniform support bound.  It makes every finite coordinate word a finite lattice
Pauli word, while the Hamel-basis condition makes every finite lattice Pauli
word a finite coordinate word.  The phase condition is independent: a
symplectic map fixes the commutator but not necessarily the chosen Pauli
multiplier.  It must also be calibrated to the phases of the actual
stabilizer group, since a label alone does not determine a $+1$ stabilizer.
The C\(^*\)-extension and GNS comparison are also distinct:
algebraic word-level surjectivity alone does not identify the weak closures in
the two representations.

Under these five assumptions the product-state GNS representation has the
spatial form
\begin{equation*}
  \pi_{\widetilde\omega}(\mathcal P_I)''
  \cong \cB(\mathcal K_P)\bar\otimes\cR_C,
\end{equation*}
where \(\mathcal K_P\) is the incomplete tensor-product Hilbert space of the pure
coordinates and \(\cR_C\) is the tracial product factor of the mixed
coordinates.  More precisely,
\begin{equation*}
  \cR_C\cong
  \begin{cases}
    M_{2^{\lvert C\rvert}}(\mathbb C),&\lvert C\rvert<\infty,\\
    \cR,&\lvert C\rvert=\infty,
  \end{cases}
  \qquad
  \dim\mathcal K_P=
  \begin{cases}
    2^{\lvert P\rvert},&\lvert P\rvert<\infty,\\
    \infty,&\lvert P\rvert=\infty.
  \end{cases}
\end{equation*}
Thus a normal tracial \emph{state} exists precisely when \(P\) is finite;
when \(P\) is infinite, the tensor product instead has the natural faithful
normal semifinite trace.  The finite type-I part of this statement is FAM-001;
the product representation, normality, and factor classification remain the
analytic work not yet expressible in the pinned Lean library.

\subsection*{Why the current manuscript proof is not imported}

The checked construction above is not a transcription of the current Overleaf
proof.  In the 2026-08-02 checkout of
\texttt{writing/overleaf\_prl\_stabilizer\_region\_algebras/main.tex}, the
following load-bearing points remain unproved or misstated:
\begin{enumerate}
\item the theorem statement at lines 1740--1752 has no short-range or adapted
finite-support symplectic-basis hypothesis, although its proof assumes one at
lines 1905--1912;
\item lines 1915--1920 declare the infinite Pauli \(*\)-isomorphism but do
not construct it, fix its phases against the stabilizer character, or prove
its C\(^*\)-extension;
\item lines 1925--1945 assert the product density matrix without proving the
precise stabilizer-membership criterion needed for all Pauli words;
\item lines 1947--1967 invoke product-GNS factorization and classification
without a normal/spatial tensor-product theorem, and the finite pure-sector
matrix size is \(2^{\lvert P\rvert}\), not the number of pure coordinates;
and
\item lines 1972--1995 do not establish equality of the two folia or a normal
isomorphism of the weak closures: the required comparison of normal preduals
and summability is left as an explicit question at lines 1992--1995.
\end{enumerate}
The formalization therefore proves the algebraic pieces that can be justified
and exposes the analytic hypotheses that must be supplied before the regional
factor statement can be promoted to a theorem.

\subsection*{Analytic proof obligations before a full theorem}

To instantiate the conditional theorem for a stabilizer region, one must prove:
\begin{enumerate}
\item an exhaustion-compatible, finite-support adapted symplectic coordinate
map; the finite-support phase correction is then given by
Proposition~\ref{prop:explicit_pauli_phase_correction} and must be calibrated
by Proposition~\ref{prop:stabilizer_character_calibration};
\item the induced quasi-local C\(^*\)-isomorphism;
\item factorization of the restricted regional state into pure and normalized-tracial
coordinate states; and
\item a spatial/normal extension identifying the bicommutants.
\end{enumerate}
Without all four, the displayed regional tensor-product isomorphism remains a
hypothesis, not a consequence of short-range entanglement.

\subsection*{Checked Lean boundary}

The pinned mathlib checkout has preliminary abstract and concrete von
Neumann-algebra definitions, but not the required tensor-product, normal-state,
hyperfiniteness, or factor-classification API.  The Lean structure
\texttt{ProductModelTraceData} records an algebra, a model
algebra, a \(*\)-equivalence, and an algebraic tracial state on the model.  Its
theorem \texttt{inducedTrace\_isAlgebraicTracialState} constructs the pullback
state.  Thus, once the analytic construction supplies the normal
tensor-product equivalence and the normal product trace, the purely algebraic
part of the asserted induced-state conclusion is already checked.  The terms
``normal,'' ``hyperfinite,'' and ``type \(II_1\)'' are not claimed by this Lean
theorem.

The tensor-matrix, dyadic-tower, and algebraic-direct-limit Lean modules
additionally verify a concrete algebraic core of the infinite qubit product:
the dyadic finite-stage traces descend to the directed limit and remain
positive on star-squares.  To obtain the hyperfinite factor, one must still
construct the relevant C\(^*\) and von Neumann completions and prove the four
analytic obligations above.

\subsection*{Sources}

The source claim appears in the Overleaf manuscript, \texttt{main.tex},
Theorem \texttt{thm:short\_range\_algebra} and the following proof discussion
(local lines 1740--1973 in the 2026-08-02 checkout).  Its explicit unresolved
markers occur at the coordinate map, product state, and tensor-product
classification steps.  For the standard model-factor facts, use Takesaki I,
Proposition IV.4.13 (represented tensor products), and Blackadar,
Sections III.3.1.4 and III.3.4.2 (tracial infinite products and
hyperfiniteness), as staged in the earlier project blueprint 012.

\end{document}
